% Fragmento integrable. Preambulo anfitrion: amsmath, amssymb, longtable, hyperref.
\section{Transport ordonné et représentation de Jacobi}
\label{hmtfg:fragmento:transporte-ordenado-y-representacion-de-jacobi}

\subsection{Transport complet depuis le continu et réalisation spectrale du profil}
\label{hmtfg:escalado:14}

Le transport ordonné du continu et la représentation spectrale du lecteur se composent avant la sélection paramétrique.

\subsubsection{Couverture de toutes les racines et transport de leur minimum}
\label{hmtfg:escalado:14.1}

La construction ordonnée du continu développée dans les travaux antérieurs produit la lecture archimédienne surjective à partir de cylindres compatibles et de leur quotient ordonné. Dans son univers déclaré \(\mathfrak G\), elle démontre l’isomorphisme de corps ordonnés complets
\begin{equation}
\iota:\mathbb R_{\rm HMT}^{\mathfrak G}
   \xrightarrow{\cong}\mathbb R^{\mathfrak G}.
\label{hmtfg:escalado:eq:37}
\end{equation}
Le nom \(\eta\) du document de référence est remplacé ici par \(\iota\), afin de le distinguer du paramètre de pondération, fixé à zéro. La seconde projection conserve l’incidence, la frontière et la mémoire sur le même antécédent ; le quotient d’évaluation et l’histoire complète conservent leurs domaines respectifs.

Construisons dans le corps HMT, avant de sélectionner une racine,
\[
\cos_{\rm H}(t)=\sum_{k=0}^{\infty}\frac{(-1)^kt^{2k}}{(2k)!},
\qquad
u_{\rm H}(x)=\frac1{12}\sum_{m=1}^{12}
\left[1-\cos_{\rm H}\!\left(\frac{2(3m-1)x}{\sqrt[\rm H]{899}}\right)\right],
\]
\begin{equation}
F_a(x)=1-a\,u_{\rm H}(x),\qquad S_n^{\rm H}(a)=F_a^{\,2^n}(0).
\label{hmtfg:escalado:eq:38}
\end{equation}
Les entiers et le deuxième moment sont ceux déjà dérivés du lecteur dodécaphasique. La racine carrée est la racine positive du corps ordonné. La série convergente réalise analytiquement ce lecteur ; elle conserve le profil et le coefficient normalisateur de la section~\ref{hmtfg:escalado:1}.

\textbf{Théorème de transport du retour.} Pour tout paramètre du domaine interne de \eqref{hmtfg:escalado:eq:38},
\begin{equation}
\iota(F_a(x))=f_{\iota(a)}(\iota(x)),\qquad
\iota(S_n^{\rm H}(a))=S_n(\iota(a)).
\label{hmtfg:escalado:eq:39}
\end{equation}
\textbf{Démonstration.} L’isomorphisme ordonné fixe les rationnels et préserve la racine positive. Il préserve les limites convergentes : chaque voisinage de rayon rationnel positif est transporté vers un autre de même rayon. Appliqué aux sommes partielles de \eqref{hmtfg:escalado:eq:38}, il transporte le cosinus et \(u_{\rm H}\). La première égalité résulte de la conservation de la somme et du produit ; la seconde, d’une récurrence sur le nombre d’itérations.

Un centre antérieur \(a_{n-1}\) étant fixé, définissons \(Z_n^{\rm H}\) par un retour nul, une période critique exacte \(2^n\) et \(a>a_{n-1}\), dans le domaine paramétrique déclaré. Soit \(Z_n\) l’ensemble défini par les mêmes conditions pour \(f_\lambda\), à droite de \(\iota(a_{n-1})\). Alors
\begin{equation}
\boxed{\iota[Z_n^{\rm H}]=Z_n,\qquad
\iota(\min Z_n^{\rm H})=\min Z_n.}
\label{hmtfg:escalado:eq:40}
\end{equation}
La seconde égalité est affirmée lorsque le minimum existe ; la première préserve également l’absence de candidats.

\textbf{Démonstration.} \eqref{hmtfg:escalado:eq:39} préserve le retour nul et les retours antérieurs qui déterminent la période exacte. L’ordre conserve la position par rapport au centre précédent. La surjectivité de \eqref{hmtfg:escalado:eq:37}, appliquée à tout paramètre racine, prouve la couverture inverse. Si \(a_*\) est le minimum, tout \(z\in Z_n\) est de la forme \(\iota(a)\), avec \(a\in Z_n^{\rm H}\), d’où \(\iota(a_*)\le z\). L’affirmation réciproque utilise \(\iota^{-1}\).

La surjectivité couvre tous les paramètres de la carte, y compris ceux dont les racines n’ont pas été isolées par un calcul fini. L’isomorphisme conserve l’univers d’interprétation déclaré et l’ordre de sélection.

\subsubsection{Représentation de Jacobi du profil dodécaphasique complet}
\label{hmtfg:escalado:14.2}

Définissons la mesure atomique du lecteur,
\begin{equation}
s_m=\frac{4(3m-1)^2}{899},\qquad
\nu_0=\frac1{12}\sum_{m=1}^{12}\delta_{s_m},\qquad
M_r=\int s^r\,d\nu_0(s).
\label{hmtfg:escalado:eq:41}
\end{equation}
Les douze nœuds sont distincts et positifs. Pour un polynôme \(p\ne0\) de degré inférieur à \(r\le12\),
\begin{equation}
\sum_{i,j=0}^{r-1}p_iM_{i+j}p_j
=\frac1{12}\sum_{m=1}^{12}p(s_m)^2>0.
\label{hmtfg:escalado:eq:42}
\end{equation}
Un polynôme de degré inférieur à douze ne s’annule pas aux douze nœuds. Au degré douze apparaît \(\prod_m(s-s_m)\), de norme nulle.

L’orthogonalisation par rapport à cette mesure produit des polynômes orthonormaux \(\psi_0,\ldots,\psi_{11}\), avec \(\psi_0=1\). Soient
\[
Q_{mj}=\psi_j(s_m)/\sqrt{12},\quad
D=\operatorname{diag}(s_m),\quad J_{12}=Q^{\mathsf T}DQ.
\]
Les relations suivantes sont satisfaites
\begin{equation}
Q^{\mathsf T}Q=I,\quad DQ=QJ_{12},\qquad
\boxed{u_0(x)=e_0^{\mathsf T}
\bigl[I-\cos(x\sqrt{J_{12}})\bigr]e_0.}
\label{hmtfg:escalado:eq:43}
\end{equation}
\textbf{Démonstration.} L’orthogonalité découle de \eqref{hmtfg:escalado:eq:42}. La multiplication par \(s\) produit une récurrence à trois termes : les éléments séparés de plus d’une position s’annulent par orthogonalité et degré. Fixer des coefficients dominants positifs rend les sous-diagonales positives. Par calcul fonctionnel,
\(\cos(x\sqrt{J_{12}})=Q^{\mathsf T}\cos(x\sqrt D)Q\).
Le vecteur \(Qe_0\) est constant, de valeur \(1/\sqrt{12}\) ; sa substitution restitue la somme de la section~\ref{hmtfg:escalado:1}.

L’article sur les grassmanniennes développe ce mécanisme pour les mesures marquées et les raffinements. Il est appliqué ici à \eqref{hmtfg:escalado:eq:41} ; le rang douze appartient à ce lecteur. Les marques généalogiques sont conservées avec le lecteur, conformément à la fidélité marquée du document de référence. \(J_{12}\) représente le profil analytique ; l’état complet conserve en outre les routes, l’orientation et la mémoire.

Le vérificateur de cette extension contrôle, avec des fractions exactes, les douze degrés de norme positive, l’annulation au degré douze, l’entrelaceur \(VT=DV\) dans la base monique et l’autoadjonction pondérée \(HT=T^{\mathsf T}H\). Il vérifie les moments de degré zéro à trente ; l’identité à tout degré découle de l’entrelaceur. Ses premiers coefficients sont
\begin{equation}
a_0=2,\qquad \beta_1=\frac{2493348}{808201}.
\label{hmtfg:escalado:eq:44}
\end{equation}

\subsubsection{Sensibilité orientée du retour}
\label{hmtfg:escalado:14.3}

Pour un retour exact de période \(p=2^n\), soient
\[
c_j=f_\lambda^j(0),\quad q_j=\partial_\lambda c_j,\quad
D_j=\prod_{k=1}^j f_\lambda'(c_k),\quad D_0=1.
\]
Avant le retour critique, les facteurs de \(D_{p-1}\) sont non nuls. La différentiation et le télescopage donnent
\begin{equation}
q_{j+1}=-u_0(c_j)+f_\lambda'(c_j)q_j,\qquad
\boxed{\frac{q_p}{D_{p-1}}
=-\sum_{j=1}^{p-1}\frac{u_0(c_j)}{D_j}.}
\label{hmtfg:escalado:eq:45}
\end{equation}
Dans le mot canonique, \(\operatorname{sign}c_j=(-1)^{v_2(j)}\). Le signe de \(f_\lambda'(c_j)\) est opposé ; par conséquent,
\begin{equation}
\operatorname{sign}D_j=(-1)^{j+\sum_{k=1}^jv_2(k)}
=(-1)^{s_2(j)},
\label{hmtfg:escalado:eq:46}
\end{equation}
en utilisant \(\sum_{k=1}^jv_2(k)=j-s_2(j)\), où \(s_2(j)\) compte les uns binaires. La sensibilité normalisée est
\begin{equation}
\mathcal T_n=-\sum_{j=1}^{2^n-1}(-1)^{s_2(j)}t_j,\qquad
t_j=\frac{u_0(c_j)}{|D_j|}>0.
\label{hmtfg:escalado:eq:47}
\end{equation}
La positivité nodale de \eqref{hmtfg:escalado:eq:42} détermine \(u_0\) ; \eqref{hmtfg:escalado:eq:47} conserve en outre les produits orbitaux et leurs orientations.

La représentation supplémentaire par moments exigerait une mesure positive \(\rho\) sur \([0,1]\) avec \(t_j=\int z^{j-1}\,d\rho\). Sous cette condition,
\begin{equation}
\mathcal T_n=\int_0^1
\frac{1-\prod_{r=0}^{n-1}(1-z^{2^r})}{z}\,d\rho(z)>0
\label{hmtfg:escalado:eq:48}
\end{equation}
pour \(\rho\ne0\), en prolongeant l’intégrande par sa valeur \(1\) en zéro. L’identité s’obtient en développant le produit binaire.

Le contrôle rationnel à la racine de période quatre précédemment certifiée obtient
\[
t_1\approx0.40425224267600139730,\quad
t_2\approx0.04925249167432646403,\quad
t_3\approx0.15740870098294833737
\]
et l’intervalle extérieur rigoureux
\begin{equation}
0.296096033367379523967764444238832345360107
<\mathcal T_2<
0.296096033367379523967764444238832345360112.
\label{hmtfg:escalado:eq:49}
\end{equation}
Il certifie également \(t_3>t_2\). Une mesure positive sur \([0,1]\) imposerait \(t_3\le t_2\) ; ce relèvement précis des poids orbitaux en moments positifs est exclu. Le signe positif de \eqref{hmtfg:escalado:eq:49} reste certifié. Ce contrôle distingue l’échec d’une preuve proposée de l’échec de la propriété que l’on cherchait à prouver.
